\pdfinfoomitdate=1
\pdftrailerid{}
\pdfsuppressptexinfo=15
\documentclass[11pt]{article}
\usepackage[T1]{fontenc}
\usepackage{lmodern}
\usepackage{amsmath,amssymb,amsthm}
\usepackage{booktabs}
\usepackage[margin=1in]{geometry}
\usepackage[hidelinks]{hyperref}
\usepackage{microtype}

\newtheorem{theorem}{Theorem}[section]
\newtheorem{lemma}[theorem]{Lemma}
\newtheorem{proposition}[theorem]{Proposition}
\newtheorem{corollary}[theorem]{Corollary}
\theoremstyle{definition}
\newtheorem{definition}[theorem]{Definition}
\newtheorem{computation}[theorem]{Verified computation}
\theoremstyle{remark}
\newtheorem{remark}[theorem]{Remark}

\newcommand{\QQ}{\mathbf Q}
\newcommand{\PP}{\mathbf P}
\newcommand{\Spec}{\operatorname{Spec}}
\newcommand{\Proj}{\operatorname{Proj}}
\newcommand{\Der}{\operatorname{Der}}
\newcommand{\Hom}{\operatorname{Hom}}
\newcommand{\MC}{\operatorname{MC}}
\newcommand{\im}{\operatorname{im}}
\newcommand{\rank}{\operatorname{rank}}
\newcommand{\ad}{\operatorname{ad}}
\newcommand{\Tor}{\operatorname{Tor}}
\newcommand{\depth}{\operatorname{depth}}
\newcommand{\pd}{\operatorname{pd}}
\newcommand{\Gr}{\operatorname{Gr}}
\newcommand{\g}{\mathfrak g}
\newcommand{\h}{\mathfrak h}

\title{Smoothing the Gr\"unbaum--Sreedharan Stanley--Reisner threefold\\
by strict differential graded Lie algebra methods}
\author{}
\date{}

\begin{document}
\maketitle

\begin{abstract}
The Stanley--Reisner threefold of the Gr\"unbaum--Sreedharan eight-vertex
triangulated three-sphere is the special fibre of a flat projective scheme
over $\QQ[[s]]$ whose geometric generic fibre is smooth.  We prove this using
only a Tate resolution $P\to A$, the strict differential graded Lie algebra
of weight-zero $S$-linear derivations of $P$, the ordinary Maurer--Cartan
equation $\partial\alpha+\tfrac12[\alpha,\alpha]=0$, elementary obstruction
theory in the complete filtered dg Lie algebra $s\,\Der_S(P,P)_0[[s]]$ with
the strict Bianchi identity, and exact rational computer algebra.  No
structure transferred to cohomology, and no operation of arity greater
than two, occurs anywhere.  The deformation-theoretic bridges are constructed
by hand rather than quoted: the certified exact family over $\QQ[s]/(s^4)$ is
lifted to an explicit square-zero perturbation of the Tate differential, and
the all-orders family is read off as the sixteen literal cubic differentials
of the degree-one Tate generators.  The generic fibre is shown to be
Cohen--Macaulay and equidimensional by a Betti-number semicontinuity argument,
after which the $280$ exact incidence certificates force smoothness.
\end{abstract}

\tableofcontents

\section{Statement, data, and conventions}\label{sec:statement}

\subsection{The special fibre}

Put $k=\QQ$ and $S=k[x_1,\ldots,x_8]$ with its standard grading.  Let
$\Delta$ be the simplicial complex on $\{1,\ldots,8\}$ with the twenty
tetrahedral facets
\[
\begin{gathered}
1234,\,1237,\,1267,\,1347,\,1567,\,2345,\,2367,\,3467,\,3456,\,4567,\\
2358,\,2368,\,3568,\,1268,\,1568,\,1248,\,2458,\,1478,\,1578,\,4578,
\end{gathered}
\]
the Gr\"unbaum--Sreedharan triangulation~\cite{GS}.  Its Stanley--Reisner
ideal is generated by the sixteen ordered monomials
\[
\begin{split}
I=(\,&m_1,\ldots,m_{16}\,)
=(\,x_6x_7x_8,\,x_4x_6x_8,\,x_3x_7x_8,\,x_3x_5x_7,\,
x_3x_4x_8,\,x_2x_7x_8,\\
&x_2x_5x_7,\,x_2x_5x_6,\,x_2x_4x_7,\,x_2x_4x_6,\,
x_1x_4x_6,\,x_1x_4x_5,\\
&x_1x_3x_8,\,x_1x_3x_6,\,x_1x_3x_5,\,x_1x_2x_5\,)\subset S.
\end{split}
\tag{1.1}\label{eq:ideal}
\]
Set $A=S/I$ and $X_0=\Proj A\subset\PP^7_k$.  For $n\geq1$ write
\[
 R_n=k[s]/(s^n),\qquad R=k[[s]],\qquad K=k((s)).
\]

\begin{theorem}[smoothing theorem]\label{thm:main}
There exist sixteen homogeneous cubics
$F_1,\ldots,F_{16}\in R[x_1,\ldots,x_8]$ with $F_i\equiv m_i\pmod s$, whose
reductions modulo $s^3$ are exactly the sixteen rows of the frozen file
\path{../referee-packet/data/universal_2jet_QQ.txt}, such that
\[
 X=\Proj\bigl(R[x_1,\ldots,x_8]/(F_1,\ldots,F_{16})\bigr)
 \longrightarrow\Spec R
\]
is flat and projective, has special fibre $X_0$, and has smooth geometric
generic fibre.  In particular $X_0$ is smoothable in characteristic zero.
\end{theorem}

This is the entire claim.  The proof does not assert smoothness of the total
space, irreducibility of the generic fibre, triviality of its canonical
bundle, or identification with any named threefold, and no such assertion is
used.

\begin{computation}[special fibre]\label{comp:special}
The exact rational replay of
\path{../referee-packet/code/verify_combinatorics.py} confirms: $\Delta$ has
$8$ vertices, $20$ facets, $97$ faces, is pure of dimension $3$, and its
inclusion-minimal nonfaces are exactly the sixteen triples in
\eqref{eq:ideal}; and for every face $\sigma$ (including $\varnothing$) the
reduced rational homology of $\operatorname{lk}_\Delta\sigma$ vanishes below
the top dimension, by exact \texttt{Fraction} row reduction of boundary
matrices.  By Reisner's criterion~\cite{Reisner}, $A$ is Cohen--Macaulay;
its Krull dimension is $\dim\Delta+1=4$, so $\depth_{S_+}A=4$.
\end{computation}

\subsection{Gradings and sign conventions}\label{subsec:conventions}

Every object below carries two gradings.

\emph{Homological degree.}  $P$ (defined in Section~\ref{sec:controller}) is
concentrated in homological degrees $\geq0$ and its differential $d$ has
degree $-1$.  A \emph{derivation of cohomological degree $p$} lowers
homological degree by $p$.

\emph{Internal weight.}  Each $x_i$ has weight $1$; every Tate generator
carries the weight of the cycle it kills ($e_i$: weight $3$; $f_j$: weight
$4$; $g_a$: weights $5$ and $6$).  Differentials and all derivations used
below preserve weight; a subscript $0$ denotes the weight-preserving part.
The parameter $s$ has weight $0$.

For homogeneous derivations $\theta,\eta$ of cohomological degrees $p,q$ set
\[
 \partial\theta=[d,\theta]=d\theta-(-1)^p\theta d,
 \qquad
 [\theta,\eta]=\theta\eta-(-1)^{pq}\eta\theta.
\tag{1.2}\label{eq:signs}
\]
The graded commutator of derivations is a derivation, $\partial^2=0$, and
$\partial$ is a degree-one derivation of the bracket, so \eqref{eq:signs}
defines strict dg Lie algebras.  These two operations, the ordinary
Maurer--Cartan equation, and no other algebraic operation, are the
deformation machinery of this paper.

\subsection{How every assertion is certified}

Each substantive statement below is one of: (i) proved here from the
definitions; (ii) a \emph{Verified computation}, meaning an exact rational
computer-algebra calculation in the frozen packets of this repository,
replayed for this document (the replay commands are in
\path{REPRODUCE.md}); (iii) a classical theorem cited with its hypotheses
checked in the text.  The classical inputs are: Reisner's
criterion~\cite{Reisner}; Tate's construction of acyclic
closures~\cite{Tate,Avramov} (re-proved inline as
Lemma~\ref{lem:tateexists}); the local flatness
criterion~\cite[Tag 00MK]{Stacks00MK} (re-proved inline in the Artinian case
used); the Auslander--Buchsbaum formula and unmixedness of Cohen--Macaulay
local rings~\cite{BH}; Hilbert--Serre dimension theory~\cite{BH}; the
valuative criterion of properness; and the finiteness of the integral
closure of a complete discrete valuation ring in a finite extension of its
fraction field~\cite[Ch.~II]{SerreLF}.  No result about structures
transferred to cohomology, of any arity, is cited or used.

\section{The Tate resolution and the two derivation algebras}
\label{sec:controller}

\subsection{The resolution}

\begin{lemma}[staged Tate resolution]\label{lem:tateexists}
There is a graded-commutative dg $S$-algebra $P$, free as a graded
commutative algebra on finitely many generators in each homological degree
$\geq1$, with weight-preserving differential $d$ of homological degree $-1$,
such that $P_0=S$, $H_0(P)=A$, and $H_i(P)=0$ for $i>0$.  It may be chosen
so that its generators in homological degrees $1,2,3$ are
\[
\begin{array}{c|c|c}
\textup{generators}&\textup{homological degree}&\textup{weight}\\ \hline
e_1,\ldots,e_{16}&1&3\\
f_1,\ldots,f_{30}&2&4\\
g_1,\ldots,g_{136}&3&5\ (16\textup{ of them}),\
6\ (120\textup{ of them}),
\end{array}
\tag{2.1}\label{eq:generators}
\]
with $d(e_i)=m_i$, with $d(f_j)$ the columns of an explicit matrix
$R_0$ over $S$, and with $d(g_a)$ the columns of an explicit matrix $Z$
written in the degree-two basis $\{f_j\}\cup\{e_ie_j\}$; the matrices are
the frozen certified data of Computation~\textup{\ref{comp:tate}} below.
\end{lemma}

\begin{proof}
This is Tate's construction~\cite{Tate,Avramov}; over $\QQ$ no divided
powers are needed and we spell it out.  Start from
$P^{(1)}=S[e_1,\ldots,e_{16}]$ with $d(e_i)=m_i$; then $H_0=A$.  Given
$P^{(q)}$ with $H_i(P^{(q)})=0$ for $1\leq i\leq q-1$, the cycle module
$Z_q(P^{(q)})$ is a finitely generated graded $S$-module, because
$P^{(q)}_q$ is a finitely generated free $S$-module and $S$ is Noetherian.
Choose finitely many weight-homogeneous cycles $z_1,\ldots,z_r$ whose
classes generate $H_q(P^{(q)})$, adjoin free (graded-commutative) generators
$w_1,\ldots,w_r$ of homological degree $q+1$ and the corresponding weights,
and set $d(w_b)=z_b$.  Every monomial involving some $w_b$ has homological
degree $\geq q+1$, so $(P^{(q+1)})_i=(P^{(q)})_i$ for $i\leq q$: hence
cycles and boundaries in degrees $<q$ are unchanged, so $H_i(P^{(q+1)})=0$
for $1\le i<q$, while in degree $q$ the cycles are unchanged and the new
boundaries $d(w_b)=z_b$ generate $H_q(P^{(q)})$, so $H_q(P^{(q+1)})=0$.  Set
$P=\bigcup_qP^{(q)}$; each element has fixed homological degree, so
$H_i(P)=H_i(P^{(i+2)})=0$ for $i\geq1$ and $H_0(P)=A$.  The differential
preserves weight because each $z_b$ was chosen weight-homogeneous.

For the specific choice \eqref{eq:generators}: the sixteen $m_i$ generate
$I$; Computation~\ref{comp:tate} certifies that the frozen $30$-column
matrix $R_0$ generates the full first syzygy module of $(m_1,\ldots,m_{16})$
and that the frozen $136$ columns, together with the boundaries of the
decomposable degree-three elements, generate the full degree-two cycle
module of the stage $S[e,f]$.  Hence the construction may take exactly these
generators in degrees $2$ and $3$ and continue with generators of degree
$\geq4$.
\end{proof}

\begin{computation}[certified low Tate stages]\label{comp:tate}
Let $D_1$, $D_2$ be the degree-one and degree-two differentials of the
algebra on \eqref{eq:generators}: $D_1$ is the row $(m_1,\ldots,m_{16})$;
$D_2$ has the $30$ columns of $R_0$ and the $120$ columns
$d(e_ie_j)=m_ie_j-m_je_i$.  Let $D_3$ be the matrix of the boundaries of the
$480+560$ decomposables $e_if_a$, $e_ie_je_k$, and let $Z$ be the
$150\times136$ matrix of the frozen $d(g_a)$.  The independent audit
\path{../audits/tate-stage/} proves, by two complete homogeneous syzygy
computations (Macaulay2) and an independent reconstruction with its own full
syzygy computations (Singular), the exact graded module equalities
\[
 D_1D_2=0,\quad D_2D_3=0,\quad D_2Z=0,\quad
 \ker D_1=\im D_2,\quad
 \ker D_2=\im[D_3\;Z].
\tag{2.2}\label{eq:tatestage}
\]
These are unbounded module identities over $S$, not finite-degree rank
tests.  The certificate is
\path{../audits/tate-stage/MATHEMATICAL_CERTIFICATE.md}.
\end{computation}

\begin{remark}[exactly what the $136$ columns certify]\label{rem:136}
Equalities \eqref{eq:tatestage} say that the $30$ generators $f_j$ kill all
homology in homological degree one of $S[e]$, and that the classes of the
$136$ cycles $d(g_a)$ span
$\ker D_2/\im D_3$, the homology in homological degree two of $S[e,f]$
after the decomposable boundaries.  \emph{No assertion that the $136$
columns span, or are independent in, $H^3$ of any cochain complex is made
anywhere in this paper, and no total dimension of any $T^3$ or $H^3$ group
is used.}  The role of \eqref{eq:tatestage} is: (i) it validates the
generator table \eqref{eq:generators} as a genuine initial Tate segment, so
that a full resolution continues with generators of homological degree
$\geq4$ only; (ii) it supplies the two cycle-lifting hypotheses used in
Lemma~\ref{lem:cyclelift} and Lemma~\ref{lem:derextend}.
\end{remark}

\subsection{Which dg Lie algebra controls what}\label{subsec:which}

Two strict dg Lie algebras act on $P$, and the proof must distinguish them.
With the conventions of \S\ref{subsec:conventions} put
\[
 \g=\Der_k(P,P)_0,
 \qquad
 \h=\Der_S(P,P)_0 .
\]
The \emph{relative} algebra $\h$ consists of derivations vanishing on
$S=P_0$; perturbing $d$ by a degree-one element of $\h$ changes nothing
about the polynomial ring or its coordinates.  It is therefore the
controller of \emph{embedded deformations in the fixed coordinates
$x_1,\ldots,x_8$}, which is what Theorem~\ref{thm:main} is about.  The
\emph{absolute} algebra $\g$ also contains degree-zero coordinate vector
fields; its $H^1$ is the abstract tangent space computed by the packages.
The two are compared, never identified:

\begin{lemma}\label{lem:relabs}
$\h^p=\g^p$ for every $p\geq1$; consequently
$Z^p(\h)=Z^p(\g)$ and $B^p(\h)\subseteq B^p(\g)$ for $p\ge 1$, and
\[
 H^1(\h)\twoheadrightarrow H^1(\g),
 \qquad
 H^q(\h)=H^q(\g)\quad(q\geq2).
\tag{2.3}\label{eq:relabscomp}
\]
\end{lemma}

\begin{proof}
A derivation of cohomological degree $p\geq1$ maps $P_0=S$ into homological
degree $-p<0$, which is zero; so it is automatically $S$-linear, giving
$\h^p=\g^p$.  For $p\geq1$ the differentials
$\partial:\h^p\to\h^{p+1}$ and $\partial:\g^p\to\g^{p+1}$ coincide, so the
cocycles agree in all degrees $\geq1$ and the coboundary spaces agree in all
degrees $\geq2$ (both are $\partial(\g^1)$).  In degree one,
$B^1(\h)=\partial(\h^0)\subseteq\partial(\g^0)=B^1(\g)$, whence the
surjection.  \emph{The surjection $H^1(\h)\to H^1(\g)$ is far from
injective}: Computation~\ref{comp:controller} gives dimensions $109$ and
$53$.
\end{proof}

\subsection{Tangent identifications}

Let $C^\bullet=\Der_k(P,A)_0$ and $C^\bullet_S=\Der_S(P,A)_0$, with
differential $\partial\theta=-(-1)^p\theta d$ (the differential of $A$ is
zero).  Since $A$ is concentrated in homological degree zero, a degree-$p$
cochain is determined by its values on the degree-$p$ generators of
\eqref{eq:generators}; explicitly, in weight zero,
\[
 C^1_S=C^1=(A_3)^{16},\quad
 C^2_S=C^2=(A_4)^{30},\quad
 C^3_S=C^3=(A_5)^{16}\oplus(A_6)^{120},
\tag{2.4}\label{eq:cochains}
\]
while $C^0=(A_1)^8$ (values on the $x_i$) and $C^0_S=0$.

\begin{lemma}[derivation quasi-isomorphisms]\label{lem:qiso}
Composition with the augmentation $P\to A$ induces quasi-isomorphisms
\[
 \g=\Der_k(P,P)_0\longrightarrow C^\bullet,
 \qquad
 \h=\Der_S(P,P)_0\longrightarrow C^\bullet_S .
\tag{2.5}\label{eq:qiso}
\]
\end{lemma}

\begin{proof}
We give the absolute case; the relative case is identical with
$\Omega_{P/S}$ in place of $\Omega_{P/k}$.  The dg $P$-module
$\Omega_{P/k}$ is free as a graded $P$-module on the basis
$\{dx_i\}\cup\{dz:\;z\ \text{a Tate generator}\}$, and
$\Der_k(P,M)=\Hom_P(\Omega_{P/k},M)$ for every dg $P$-module $M$.  Filter
$\Omega_{P/k}$ by the sub-dg-modules $F_q$ generated by the $dx_i$ and the
$dz$ with $z$ of homological degree $\leq q$.  For a generator $z$ of degree
$q$, $d(dz)=d(d_Pz)$ is the universal derivative of an element of the
subalgebra on earlier generators, hence lies in $F_{q-1}$; so each $F_q$ is
$d$-stable and $F_q/F_{q-1}$ is a free $P$-module on classes with zero
induced differential.  Let $Q=\ker(P\to A)$; it is acyclic, because $P\to A$
is a surjective quasi-isomorphism.  For each $q$,
$\Hom_P(F_q/F_{q-1},Q)$ is a product of shifted copies of $Q$, hence
acyclic (homology commutes with direct products of complexes of vector
spaces), and by induction each $\Hom_P(F_q,Q)$ is acyclic (a degreewise-split
extension of acyclic complexes).  The restrictions
$\Hom_P(F_q,Q)\to\Hom_P(F_{q-1},Q)$ are surjective because
$F_{q-1}\subset F_q$ splits as graded modules; hence in the inverse limit
$\Hom_P(\Omega_{P/k},Q)=\varprojlim_q\Hom_P(F_q,Q)$ the
$\varprojlim^1$-terms of the Milnor sequence vanish and the limit is
acyclic.  Finally $0\to\Hom(\Omega,Q)\to\Hom(\Omega,P)\to\Hom(\Omega,A)\to0$
is exact because $\Omega$ is graded-free, so the second map is a
quasi-isomorphism.  All modules split as direct products over internal
weights, homology commutes with those products, and therefore the
weight-zero components in \eqref{eq:qiso} are quasi-isomorphisms as well.
\end{proof}

\begin{lemma}[cohomology in low degrees]\label{lem:lowcoh}
With identifications \eqref{eq:cochains}:
\begin{enumerate}
\item $H^1(\h)=H^1(C_S)=\Hom_S(I,A)_0$, the space of weight-zero
$S$-module maps $I\to A$, via $\theta\mapsto(m_i\mapsto\theta(e_i))$.
\item $H^1(\g)=H^1(C)=\Hom_S(I,A)_0/\im(\textup{Jac})$, where
$\textup{Jac}$ sends a weight-zero vector field
$\xi=\sum\xi_i\partial/\partial x_i$ to $(m_i\mapsto\xi(m_i))$.  This is
the classical embedded-to-abstract quotient defining $T^1_0(A)$.
\item $H^2(\h)=H^2(\g)=T^2_0(A)$ and $H^3(\h)=H^3(\g)$, computed in
$C^\bullet$ as $\ker\delta^2/\im\delta^1$ and $\ker\delta^3/\im\delta^2$
respectively, where $\delta^1$ is (the transpose of) $R_0$ read in $A$ and
$\delta^2$ is the linear $f$-part of the $d(g_a)$ read in $A$.
\end{enumerate}
\end{lemma}

\begin{proof}
(1) $C^0_S=0$, so $H^1(C_S)=\ker(\delta^1)$; a degree-one cochain with
values $v_i=\theta(e_i)\in A_3$ is closed iff
$\theta(d f_j)=\sum_i (R_0)_{ij}v_i=0$ in $A$ for all $j$.  Because $R_0$
generates \emph{all} first syzygies of $(m_1,\ldots,m_{16})$
(Computation~\ref{comp:tate}), this is exactly the condition that
$m_i\mapsto v_i$ is a well-defined $S$-map $I\to A$.
(2) The absolute complex has the extra term $C^0=(A_1)^8$ with
$\partial\xi$ acting on $e_i$ by $\mp\xi(m_i)$; the image is the Jacobian
image, and the sign does not change it.
(3) For $q\geq2$ the two cochain complexes agree from degree one onward
except in degree zero and one, and coboundaries in degree $\geq2$ come from
degree $\geq1$, where the complexes agree; combine with
Lemma~\ref{lem:qiso}.  A degree-two derivation kills the $e_i$ (their
values would land in negative homological degree of $A$) and hence kills the
decomposables $e_ie_j$; so $\delta^2\eta(g_a)=\pm\eta(dg_a)$ sees only the
linear $f$-part of $dg_a$.  The identification of $H^2(C)$ with the
package's $T^2_0(A)$, at the level of the specific bases used, is part of
Computation~\ref{comp:bracket} below.
\end{proof}

\begin{computation}[controller dimensions]\label{comp:controller}
The exact script \path{verify_fable_dgla.m2} in this directory computes
\[
 \dim\Hom_S(I,A)_0=109,\qquad
 \dim\im(\textup{Jac})_0=56,\qquad
 \dim T^1_0(A)=53=109-56,
\]
and $\dim T^2_0(A)=27$.  The certificate is
\path{verification/fable_dgla_QQ.txt}.  Thus $H^1(\h)$ has dimension $109$
and $H^1(\g)$ dimension $53$; the proof never conflates them.
\end{computation}

\section{Maurer--Cartan elements give embedded flat graded families}
\label{sec:mc}

For $n\geq1$ set $\h_n=\h\otimes_k(s)/(s^n)$, a nilpotent dg Lie algebra
over $R_n$, and let
\[
 \widehat\h=s\,\h[[s]]=\prod_{j\geq1}s^j\h,
 \qquad F^r\widehat\h=s^r\h[[s]] ,
\tag{3.1}\label{eq:filtration}
\]
a dg Lie algebra with a complete, separated filtration for which brackets
add filtration levels.  The Maurer--Cartan sets are
\[
 \MC(\h_n)=\Bigl\{\alpha\in\h_n^1:\ \partial\alpha+\tfrac12[\alpha,\alpha]=0
 \Bigr\},
 \qquad\text{and likewise } \MC(\widehat\h).
\]

\begin{proposition}[realization]\label{prop:realization}
Let $\alpha\in\MC(\h_n)$.  Then $D=d+\alpha$ is an $R_n[x]$-linear,
weight-preserving, square-zero derivation of $P\otimes_kR_n$ of homological
degree $-1$, and:
\begin{enumerate}
\item $H_i(P\otimes R_n,D)=0$ for $i>0$;
\item $B_n:=H_0(P\otimes R_n,D)=R_n[x_1,\ldots,x_8]/\bigl(D(e_1),\ldots,
D(e_{16})\bigr)$, and each $D(e_i)$ is a homogeneous cubic congruent to
$m_i$ modulo $s$;
\item every weight component $(B_n)_j$ is a free $R_n$-module of rank
$\dim_k A_j$; in particular $B_n$ is a flat graded embedded deformation of
$A$ over $R_n$ in the fixed coordinates.
\end{enumerate}
The same statements hold over $R$ for $\alpha\in\MC(\widehat\h)$, with
\textup{(3)} replaced by: each weight component of
$B=H_0(P\otimes R,D)$ satisfies $B/s^nB=B_n$ for all $n$ with each
$(B_n)_j$ free.
\end{proposition}

\begin{proof}
For a degree-one derivation $\alpha$ one has
$(d+\alpha)^2=\tfrac12[d+\alpha,d+\alpha]
=\partial\alpha+\tfrac12[\alpha,\alpha]$ (using $d^2=0$), so the
Maurer--Cartan equation is literally $D^2=0$.  $S$-linearity of $\alpha$ and
of $d$ gives $R_n[x]$-linearity of $D$; weight preservation holds because
$\alpha\in\h_0$ by definition.

(1) Induct on $n$.  The sub-complex $s^{n-1}(P\otimes R_n)$ is $D$-stable
and, since $\alpha\cdot s^{n-1}\equiv0$, is isomorphic to $(P,d)$; the
quotient is $(P\otimes R_{n-1},D\bmod s^{n-1})$.  The long exact homology
sequence and the inductive hypothesis (base $n=1$: $(P,d)$ itself) give
$H_i=0$ for $i>0$ and the exact rows
\[
 0\to A\to H_0(P\otimes R_n,D)\to H_0(P\otimes R_{n-1},D)\to0 .
\tag{3.2}\label{eq:h0les}
\]

(2) $(P\otimes R_n)_0=R_n[x]$ and $(P\otimes R_n)_1$ is the free
$R_n[x]$-module on $e_1,\ldots,e_{16}$; by $R_n[x]$-linearity the image of
$D$ in degree zero is the ideal generated by the $D(e_i)$.  Each
$D(e_i)=m_i+\alpha(e_i)$ has weight $3$, i.e.\ is a cubic in $x$ with
coefficients in $(s)\subset R_n$ beyond its leading term $m_i$.

(3) Fix a weight $j$.  Each $(P\otimes R_n)_{i,j}$ is a finite free
$R_n$-module: there are finitely many Tate generators in each homological
degree (Lemma~\ref{lem:tateexists}), each of weight $\geq3$, so a fixed
pair $(i,j)$ admits only finitely many generator monomials, each with a
finite-dimensional space of polynomial coefficients.  By (1) and (2),
\[
 \cdots\to(P\otimes R_n)_{1,j}\xrightarrow{\;D\;}
 (P\otimes R_n)_{0,j}\to(B_n)_j\to0
\]
is a resolution of $(B_n)_j$ by finite free $R_n$-modules.  Applying
$-\otimes_{R_n}k$ returns the complex $(P_{\ast,j},d)$, because
$D\equiv d\pmod s$; its positive homology vanishes.  Hence
$\Tor^{R_n}_1(k,(B_n)_j)=H_1(P_{\ast,j},d)=0$.  Now use the local flatness
criterion for the finite module $(B_n)_j$ over the Artinian local ring
$R_n$: choose a surjection $G\to(B_n)_j$ from a finite free module with
$G\otimes k\cong(B_n)_j\otimes k$ and kernel $Z\subseteq(s)G$; vanishing of
$\Tor_1$ makes $0\to Z\otimes k\to G\otimes k\to(B_n)_j\otimes k\to0$ exact,
so $Z\otimes k=0$ and $Z=0$ by Nakayama.  (This is the Artinian case of
\cite[Tag 00MK]{Stacks00MK}.)  Thus $(B_n)_j$ is free; its rank equals
$\dim_k(B_n)_j\otimes_{R_n}k$, and since $H_0$ is a cokernel and hence
commutes with the base change $R_n\to k$, this is $\dim_kA_j$.

Over $R$: reduction of $D$ modulo $s^n$ is the realization of
$\alpha\bmod s^n\in\MC(\h_n)$, and
$H_0(P\otimes R,D)/s^n\cong H_0(P\otimes R_n,D_n)$ in each weight by
\eqref{eq:h0les} and right-exactness; the displayed statements follow.
\end{proof}

\begin{remark}\label{rem:noconverse}
No converse (``every deformation comes from a Maurer--Cartan element'') and
no equivalence of deformation functors is invoked anywhere in this paper.
The one deformation we must convert into a Maurer--Cartan element---the
certified family over $R_4$---is converted by the explicit construction of
Section~\ref{sec:start}.  This is why no deformation-functor formalism, in
dg Lie or any other language, appears among the dependencies.
\end{remark}

The extension tool for that construction is:

\begin{lemma}[filtered cycle lifting]\label{lem:cyclelift}
Let $Q\subseteq P$ be a $d$-stable graded subalgebra containing $S$ and all
Tate generators of homological degree $<q$, where $q\geq3$, and suppose
$H_{q-2}(Q,d)=0$.  Let $D=d+\beta$ be a square-zero weight-preserving
differential on $Q\otimes R_n$ with $\beta$ $S$-linear of positive
$s$-order.  Then every weight-homogeneous $d$-cycle
$c_0\in Q_{q-1}$ admits a weight-homogeneous $D$-cycle
$c\in Q_{q-1}\otimes R_n$ with $c\equiv c_0\pmod s$.
\end{lemma}

\begin{proof}
Build $c^{(r)}=c_0+su_1+\cdots+s^{r-1}u_{r-1}$ with
$Dc^{(r)}\equiv0\pmod{s^r}$; the case $r=1$ is $dc_0=0$.  Write
$Dc^{(r)}\equiv s^r\rho\pmod{s^{r+1}}$ with $\rho\in Q_{q-2}$.  Applying $D$
and using $D^2=0$ gives $s^rd\rho\equiv0\pmod{s^{r+1}}$, so $d\rho=0$.
Since $H_{q-2}(Q)=0$, choose $u_r\in Q_{q-1}$ with $du_r=-\rho$, homogeneous
of the required weight (both $d$ and $D$ preserve weight, so $\rho$ is
weight-homogeneous and the certified graded syzygy computations allow a
homogeneous primitive).  Then $c^{(r+1)}$ improves the congruence; after
$n-1$ steps $c=c^{(n)}$ is a $D$-cycle.
\end{proof}

\section{The certified $R_4$ family is a genuine truncated Maurer--Cartan
element}\label{sec:start}

\subsection{The certified matrices}

Fix the rational tangent vector, in the ordered $53$-column basis
$T^1$ computed by \texttt{CT\^{}1(0,F)}:
\[
\begin{split}
y=(\,&1,3,1,2,6,1,1,1,1,1,21,28,1,5,10,55,66,15,1,18,2,\\
&35,42,1,4,10,1,25,12,30,5,33,44,1,6,3,14,7,1,9,1,6,\\
&12,20,1,4,22,1,15,1,1,8,11\,).
\end{split}
\tag{4.1}\label{eq:y}
\]
No property of the provenance of \eqref{eq:y} is used.

\begin{computation}[the starting family]\label{comp:start}
Write $F_0,\ldots,F_3$ and $Q_0,\ldots,Q_3$ for the order-$0$ to order-$3$
family and relation-lift matrices returned by the pinned
\texttt{versalDeformation} run (\texttt{HighestOrder=>3},
\texttt{SmartLift=>false}), and specialize the parameters along $t=sy$.  Put
$F^{[3]}=\sum_{i\le3}F_i$ and $Q^{[3]}=\sum_{i\le3}Q_i$ over
$R_4[x_1,\ldots,x_8]$.  The exact scripts
\path{../completion/verify_starting_jet.m2} and
\path{verify_fable_dgla.m2} verify:
\begin{enumerate}
\item the quadratic and cubic base equations vanish at $y$, and
\[
 F^{[3]}\,Q^{[3]}=0 \quad\text{over } R_4[x];
\tag{4.2}\label{eq:fq}
\]
\item $Q_0$ is the \emph{complete} $30$-column first-syzygy matrix of the
row $(m_1,\ldots,m_{16})$;
\item \emph{weight homogeneity}: every entry of $F^{[3]}$ is homogeneous of
$x$-degree $3$ and every entry of $Q^{[3]}$ is homogeneous of $x$-degree
$1$ (new check in \path{verify_fable_dgla.m2});
\item the first-order coefficient of $F^{[3]}$ equals $T^1y$ entrywise;
\item the reduction of $F^{[3]}$ modulo $s^3$ regenerates, byte for byte,
the frozen sixteen-row two-jet file with SHA-256
\[
\texttt{40e64e61674b6a4e61f1ea6822dc79327bf4ba397285f84ed8738c4c18cd1795}.
\]
\end{enumerate}
Certificates: \path{../completion/verification/starting_jet_QQ.txt},
\path{verification/fable_dgla_QQ.txt}.
\end{computation}

\subsection{Lifting the matrices to a square-zero Tate differential}

The matrices of Computation~\ref{comp:start} are \emph{not} treated as
cohomology classes.  They are used as the values of an explicit derivation.

The Tate model's relation matrix $R_0$ (Lemma~\ref{lem:tateexists}) and the
package syzygy matrix $Q_0$ generate the same module; the frozen script
\path{../referee-packet/code/verify_aq_bracket.m2} computes (and every
replay reverifies) a constant $30\times30$ matrix $U$ with
\[
 Q_0U=R_0,\qquad \det U=1 .
\tag{4.3}\label{eq:U}
\]

\begin{proposition}[the truncated Maurer--Cartan element]\label{prop:start}
Define a degree $-1$, weight-preserving, $R_4[x]$-linear derivation $D$ on
$P\otimes R_4$ by
\[
 D(e_i)=F^{[3]}_i,
 \qquad
 D(f_j)=\sum_{i=1}^{16}\bigl(Q^{[3]}U\bigr)_{ij}\,e_i,
\tag{4.4}\label{eq:Ddef}
\]
and on the Tate generators of homological degree $\geq3$ by
Lemma~\textup{\ref{lem:cyclelift}}, as detailed in the proof.  Then $D^2=0$,
$D\equiv d\pmod s$, and
\[
 \bar\alpha:=D-d\ \in\ \MC\bigl(\h\otimes(s)/(s^4)\bigr).
\tag{4.5}\label{eq:alphabar}
\]
Moreover $H_0(P\otimes R_4,D)=R_4[x]/(F^{[3]})=:A_4$, a flat graded embedded
deformation whose reduction modulo $s^3$ is the printed two-jet, and the
linear coefficient $a$ of $\bar\alpha$ is a weight-zero cocycle in $\h^1$
with $a(e_i)=(T^1y)_i$.
\end{proposition}

\begin{proof}
On the generators \eqref{eq:Ddef}: $D^2(e_i)=D(F^{[3]}_i)=0$ by
$R_4[x]$-linearity, and
\[
 D^2(f_j)=\sum_i(Q^{[3]}U)_{ij}F^{[3]}_i
 =\bigl(F^{[3]}Q^{[3]}U\bigr)_j=0
\]
by \eqref{eq:fq}.  Weight preservation on $e,f$ is
Computation~\ref{comp:start}(3).  Reduction modulo $s$: $F^{[3]}\equiv
(m_i)$ and $Q^{[3]}U\equiv Q_0U=R_0$, so $D\equiv d$ on $e$ and $f$.

Extend over the remaining generators by induction on homological degree.
For a generator $z$ of degree $q\geq3$, the element $dz$ is a
weight-homogeneous $d$-cycle of degree $q-1$ in the subalgebra $P_{<q}$
on the earlier generators.  The hypothesis $H_{q-2}(P_{<q},d)=0$ of
Lemma~\ref{lem:cyclelift} holds: for $q=3$ it is
$\ker D_1=\im D_2$ and for $q=4$ it is $\ker D_2=\im[D_3\,Z]$, both
certified in \eqref{eq:tatestage}; for $q\geq5$ it holds because
$P_{<q}=P^{(q-1)}$ is the Tate stage with $H_i=0$ for $1\leq i\leq q-2$
(Lemma~\ref{lem:tateexists}).  Also $D$ restricted to $P_{<q}\otimes R_4$
squares to zero by the induction (its square, an even derivation, vanishes
on generators).  Lemma~\ref{lem:cyclelift} produces a weight-homogeneous
$D$-cycle $c_z\equiv dz\pmod s$; set $D(z)=c_z$.  Then $D^2(z)=D(c_z)=0$,
and $D\equiv d\pmod s$ on $z$.  The extension of a value set from
generators to all of $P\otimes R_4$ by the signed Leibniz rule is a
derivation, and the square of the odd derivation $D$ is the derivation
$\tfrac12[D,D]$, so $D^2=0$ everywhere.

Consequently $\bar\alpha=D-d$ is $S$-linear (both $D$ and $d$ vanish on
$S\otimes R_4$, being $R_4[x]$-linear by definition),
weight-zero, and of degree one with coefficients in $(s)/(s^4)$, and
$D^2=0$ is \eqref{eq:alphabar}.  Proposition~\ref{prop:realization}
identifies $H_0=R_4[x]/(D(e_i))=A_4$ and gives flatness; the reduction of
$F^{[3]}$ modulo $s^3$ is the printed two-jet by
Computation~\ref{comp:start}(5).  The Maurer--Cartan equation modulo $s^2$
reads $\partial a=0$, and $a(e_i)$ is the $s$-coefficient of
$F^{[3]}_i$, which is $(T^1y)_i$ by Computation~\ref{comp:start}(4).
\end{proof}

\section{The exact sequence at $H^2$}\label{sec:exact}

\subsection{The computed derivations and their global closure}

The frozen script \path{../referee-packet/code/verify_aq_bracket.m2}
constructs, by exact linear algebra over $\QQ$:
a degree-one derivation $\theta$ with $\theta(e_i)=(T^1y)_i$, solved to be
closed on all generators \eqref{eq:generators}; degree-one derivations
$\psi_1,\ldots,\psi_{53}$ with $\psi_t(e_i)=(T^1)_{it}$, solved to be
closed through the $f_j$; and degree-two derivations
$\eta_1,\ldots,\eta_{27}$ whose $f$-values are the transport by
$U^{\mathsf T}$ of the package $T^2$ basis, checked to be closed on all
$g_a$.

\begin{lemma}[global extension of the computed derivations]
\label{lem:derextend}
$\theta$ and each $\eta_j$ extend, without changing their values on
\eqref{eq:generators}, to closed derivations of the full Tate resolution
$P$; each $\psi_t$ extends without changing its values on the $e_i,f_j$.
\end{lemma}

\begin{proof}
Let $\vartheta$ be a derivation of cohomological degree $p\in\{1,2\}$,
closed on all generators of homological degree $<q$.  Closure on a
generator $z$ of degree $q$ demands
$d(\vartheta z)=(-1)^p\vartheta(dz)$; the right side is a $d$-cycle (apply
closure on earlier generators to $d(dz)=0$) of homological degree
$q-p-1$.  For the first open case of $\psi_t$, namely $p=1$, $q=3$, that
degree is $1$ and $\ker D_1=\im D_2$ \eqref{eq:tatestage} supplies a
weight-homogeneous primitive in $P_{<3}$.  For all cases with $q\geq4$ the
degree $q-p-1$ is positive and lies in the stage $P^{(q-1)}$, which is
acyclic there (Lemma~\ref{lem:tateexists}), so a primitive among earlier
generators exists; choose its component of the required weight.  Induction
over the well-ordered generators finishes.  The values already checked are
never altered.
\end{proof}

\begin{corollary}\label{cor:h3}
The $27$ commutators $[\theta,\eta_j]=\theta\eta_j-\eta_j\theta$ are closed
degree-three derivations of $P$; their images in $C^3$ lie in
$\ker\delta^3$.  Hence, on their span, the projection
\[
 H^3(\g)=\ker\delta^3/\im\delta^2
 \hookrightarrow C^3/\im\delta^2
\tag{5.1}\label{eq:h3inject}
\]
computes rank in genuine $H^3(\g)$.  Likewise the $53$ commutators
$[\theta,\psi_t]=\theta\psi_t+\psi_t\theta$ are closed degree-two
derivations, representing classes of $H^2(\g)$.
\end{corollary}

This is the rigorous content behind the ``$136$ cycles'' step: the $136$
generators enter only through the certified completeness
\eqref{eq:tatestage}, which (i) makes the displayed $g$-stage the correct
domain for degree-three cochain values, and (ii) guarantees the closure
solutions and extensions above.  Nothing about a basis of $H^3$ is claimed.

\begin{computation}[bracket data]\label{comp:bracket}
The frozen script, replayed exactly, proves over $\QQ$:
\begin{enumerate}
\item the transported $27$ columns are $\partial$-closed, independent
modulo $\im\delta^1$, and no weight-zero cocycle in $C^2$ remains outside
their span plus $\im\delta^1$; hence $[\eta_1],\ldots,[\eta_{27}]$ is a
basis of the $27$-dimensional $H^2(\g)$;
\item the matrix of $\bigl([\theta,\eta_j]\bigr)_{j\le 27}$ in
$C^3/\im\delta^2$ has rank $12$;
\item the matrix of $\bigl([\theta,\psi_t]\bigr)_{t\le53}$ in
$C^2/\im\delta^1$ equals $-\,J\,Dq_y$, where $J$ is the (injective modulo
boundaries) matrix of the $27$ classes and $Dq_y$ is the Jacobian at $y$ of
the $27$ quadratic package equations; $\rank Dq_y=15$;
\item the composite of the matrices in \textup{(2)} and \textup{(3)} is
zero, and the column $\bigl([\theta,\psi_t]\bigr)y$ vanishes.
\end{enumerate}
Certificate: \path{../referee-packet/verification/aq_bracket_QQ.txt}
(replayed lines: \texttt{rank commutator columns modulo delta2=12},
\texttt{rank Dq\_y=15}, \texttt{primary plus Dq rank=0},
\texttt{ad\_y composed primary-bracket rank=0},
\texttt{Jacobi composite rank=0}).
\end{computation}

\subsection{From the script's cocycle to the actual linear coefficient}

Let $a$ be the linear coefficient of the constructed element
\eqref{eq:alphabar} and $\theta$ the script's cocycle.

\begin{lemma}[chain-level bridge]\label{lem:bridge}
$a-\theta=\partial\xi$ for some $\xi\in\g^0$.  Consequently, for every
closed $z\in\h^p$ with $p\geq1$,
\[
 [a,z]-[\theta,z]=\partial[\xi,z],
 \qquad [\xi,z]\in\h^p,
\]
so $\ad_a$ and $\ad_\theta$ induce the same maps
$H^p(\h)\to H^{p+1}(\h)$ for all $p\geq1$.
\end{lemma}

\begin{proof}
Both $a$ and $\theta$ are closed ($\partial a=0$ from the Maurer--Cartan
equation modulo $s^2$; $\theta$ by construction), and both take the value
$(T^1y)_i\in S_3$ on $e_i$, \emph{literally as polynomials}
(Proposition~\ref{prop:start} and the definition of $\theta$).  A
degree-one element of $C^1=\Der_k(P,A)_0$ is determined by its $e$-values;
hence the image of $a-\theta$ in $C^1$ is the zero cochain.  By
Lemma~\ref{lem:qiso} the map $H^1(\g)\to H^1(C)$ is injective, so
$[a-\theta]=0$ in $H^1(\g)$, i.e.\ $a-\theta=\partial\xi$ with
$\xi\in\g^0$.  For closed $z$ of degree $p\geq1$, the derivation property
of $\partial$ gives
$\partial[\xi,z]=[\partial\xi,z]+(-1)^0[\xi,\partial z]=[a-\theta,z]$,
and $[\xi,z]$ has cohomological degree $p\geq1$, hence lies in
$\h^p=\g^p$ (Lemma~\ref{lem:relabs}).
\end{proof}

\subsection{Exactness}

\begin{proposition}[exactness at $H^2$]\label{prop:exact}
With $a$ the linear coefficient of \eqref{eq:alphabar},
\[
 H^1(\h)\xrightarrow{\;\ad_a\;}H^2(\h)\xrightarrow{\;\ad_a\;}H^3(\h)
\tag{5.2}\label{eq:sequence}
\]
is exact at the $27$-dimensional middle term; the first map has image of
dimension $15$ and the second has rank $12$.
\end{proposition}

\begin{proof}
By Lemma~\ref{lem:bridge} it suffices to prove the statement for
$\ad_\theta$.  We first prove it for the absolute algebra $\g$ and then
transfer.

\emph{Absolute exactness.}
By Computation~\ref{comp:bracket}(1) and Lemma~\ref{lem:lowcoh}(3), the
$[\eta_j]$ are a basis of $H^2(\g)$.  By Corollary~\ref{cor:h3} and
\eqref{eq:h3inject}, the rank $12$ of Computation~\ref{comp:bracket}(2) is
the rank of $\ad_\theta:H^2(\g)\to H^3(\g)$; note that this uses the
injection \eqref{eq:h3inject} only on already-closed columns, so no
dimension of $H^3(\g)$ is needed or claimed.  By
Computation~\ref{comp:bracket}(3), the image of
$\ad_\theta:H^1(\g)\to H^2(\g)$ is the column space of $J\,Dq_y$; since the
$27$ columns of $J$ are independent modulo boundaries, this image has
dimension $\rank Dq_y=15$.  By Computation~\ref{comp:bracket}(4) the
composite vanishes on the computed matrices, i.e.\
$\im\ad_\theta\subseteq\ker\ad_\theta$ at $H^2$.  Finally
\[
 \dim\ker\bigl(\ad_\theta:H^2\to H^3\bigr)=27-12=15
 =\dim\im\bigl(\ad_\theta:H^1\to H^2\bigr),
\]
so the inclusion is an equality.  \emph{The exactness is thus obtained from
the two independently computed ranks together with the directly verified
zero composite---not from the single rank of $Dq_y$.}

\emph{Transfer to $\h$.}  By Lemma~\ref{lem:relabs},
$H^2(\h)=H^2(\g)$ and $H^3(\h)=H^3(\g)$, and the second map in
\eqref{eq:sequence} is the absolute one; so
$\ker(\ad_\theta:H^2(\h)\to H^3(\h))$ has dimension $15$.  For the image of
the first map: the surjection $H^1(\h)\to H^1(\g)$ of
Lemma~\ref{lem:relabs} intertwines the two maps $\ad_\theta$, so
$\ad_\theta H^1(\h)\supseteq$ the image computed from any absolute class
represented by a relative cocycle---and every absolute class is so
represented since $Z^1(\h)=Z^1(\g)$.  Conversely if a relative class dies
absolutely, i.e.\ is represented by $\partial\xi$ with $\xi\in\g^0$, then
\[
 [\theta,\partial\xi]=-\partial[\theta,\xi],
\]
because $\partial\theta=0$ and $\partial$ is a degree-one derivation of the
bracket, and $[\theta,\xi]\in\g^1=\h^1$; so such classes contribute nothing
to the image in $H^2(\h)$.  Hence
$\im(\ad_\theta:H^1(\h)\to H^2(\h))
=\im(\ad_\theta:H^1(\g)\to H^2(\g))$, of dimension $15$, and
\eqref{eq:sequence} is exact at $H^2$.
\end{proof}

\begin{remark}
The $56$ extra relative tangent classes of
Computation~\ref{comp:controller} are precisely the coordinate directions
that become absolute boundaries; the proof above shows they are invisible
to the obstruction calculus.  This is why the $109/53$ discrepancy---which
would be fatal if the two $H^1$'s were confused---is harmless when tracked
correctly.
\end{remark}

\section{The strict all-orders induction}\label{sec:induction}

The abstract lemma is stated and proved in full, with all signs, in the
companion note \path{DG_LIE_INDUCTION.tex}; we reproduce it here so that
this document is standalone.

For $\gamma\in\widehat\h^1$ define the curvature
\[
 \mathcal F(\gamma)=\partial\gamma+\tfrac12[\gamma,\gamma].
\tag{6.1}\label{eq:curv}
\]

\begin{lemma}[strict Bianchi identity]\label{lem:bianchi}
$\partial\mathcal F(\gamma)+[\gamma,\mathcal F(\gamma)]=0$.
\end{lemma}

\begin{proof}
Since $\partial$ is a degree-one derivation of the bracket and
$|\gamma|=1$:
$\partial[\gamma,\gamma]=[\partial\gamma,\gamma]-[\gamma,\partial\gamma]
=-2[\gamma,\partial\gamma]$, using
$[\partial\gamma,\gamma]=-(-1)^{2\cdot1}[\gamma,\partial\gamma]$.  Hence
$\partial\mathcal F(\gamma)=-[\gamma,\partial\gamma]$.  On the other hand
$[\gamma,\mathcal F(\gamma)]=[\gamma,\partial\gamma]
+\tfrac12[\gamma,[\gamma,\gamma]]$, and the graded Jacobi identity applied
to three copies of the odd element $\gamma$ gives
$3[\gamma,[\gamma,\gamma]]=0$, hence $[\gamma,[\gamma,\gamma]]=0$ over
$\QQ$.  Adding the two displays gives the identity.
\end{proof}

\begin{lemma}[fixed-two-jet extension]\label{lem:induction}
Let $\bar\alpha\in\MC(\h\otimes(s)/(s^4))$ have linear coefficient $a$, and
suppose \eqref{eq:sequence} is exact at $H^2(\h)$.  Then there is
\[
 \alpha_\infty\in\MC(\widehat\h)
 \qquad\text{with}\qquad
 \alpha_\infty\equiv\bar\alpha\pmod{s^3}.
\tag{6.2}\label{eq:allorders}
\]
\end{lemma}

\begin{proof}
Choose any lift $\gamma^{(4)}\in\widehat\h^1$ of $\bar\alpha$ (for
instance with no coefficients beyond $s^3$); then
$\mathcal F(\gamma^{(4)})\in F^4\widehat\h$ because the Maurer--Cartan
equation holds modulo $s^4$.  Inductively assume, for $N\geq4$, an element
$\gamma^{(N)}=sa+O(s^2)$ with
\[
 \mathcal F(\gamma^{(N)})=s^Nr_N+s^{N+1}r_{N+1}+O(s^{N+2}),
 \qquad r_i\in\h^2 .
\tag{6.3}\label{eq:indhyp}
\]

\emph{Step 1: the leading obstruction is a cocycle.}  The coefficient of
$s^N$ in the Bianchi identity of Lemma~\ref{lem:bianchi} is
$\partial r_N=0$ (all bracket terms pair a curvature coefficient, of order
$\geq N$, with a coefficient of $\gamma^{(N)}$, of order $\geq1$).

\emph{Step 2: its class lies in $\ker\ad_a$.}  The coefficient of
$s^{N+1}$ in the Bianchi identity is
\[
 \partial r_{N+1}+[a,r_N]=0,
\]
so $[a,r_N]$ is exact and $[r_N]\in\ker(\ad_a:H^2(\h)\to H^3(\h))$.

\emph{Step 3: kill it by a closed degree-one correction plus a boundary.}
By exactness (Proposition~\ref{prop:exact}) there is a \emph{closed}
$v\in Z^1(\h)$ with $[r_N]=-[[a,v]]$ in $H^2(\h)$; choose
$w\in\h^1$ with
\[
 r_N+[a,v]=\partial w .
\tag{6.4}\label{eq:vw}
\]
Set $\gamma^{(N+1)}=\gamma^{(N)}+\delta$ with
$\delta=s^{N-1}v-s^Nw$.  The exact change formula (valid in any dg Lie
algebra, for $|\delta|=1$) is
\[
 \mathcal F(\gamma+\delta)-\mathcal F(\gamma)
 =\partial\delta+[\gamma,\delta]+\tfrac12[\delta,\delta].
\tag{6.5}\label{eq:change}
\]
Here $\partial\delta=-s^N\partial w$ (as $\partial v=0$);
$[\gamma^{(N)},\delta]=s^N[a,v]+O(s^{N+1})$, because
$\gamma^{(N)}-sa=O(s^2)$ and $[\gamma^{(N)},s^Nw]=O(s^{N+1})$; and
$\tfrac12[\delta,\delta]=O(s^{2N-2})$ with $2N-2\geq N+2$ for $N\geq4$.  So
the coefficient of $s^N$ changes by $[a,v]-\partial w=-r_N$ by
\eqref{eq:vw}, no lower coefficient changes, and
$\mathcal F(\gamma^{(N+1)})=O(s^{N+1})$, restoring \eqref{eq:indhyp} with
$N+1$.

\emph{Step 4: convergence and the fixed two-jet.}  The step indexed by $N$
changes $\gamma$ only in $s$-orders $\geq N-1$; hence for each fixed $j$
the coefficient of $s^j$ stabilizes after step $N=j+1$, and by completeness
of the filtration \eqref{eq:filtration} the sequence converges
coefficientwise to some $\alpha_\infty\in\widehat\h^1$.  For each $j$, the
$s^j$-coefficient of $\mathcal F(\alpha_\infty)$ is a finite expression in
the first $j$ coefficients of $\alpha_\infty$, which agree with those of
$\gamma^{(N)}$ for $N>j+1$; since $\mathcal F(\gamma^{(N)})\in
F^N\widehat\h$, that coefficient is zero.  Separatedness
($\bigcap_NF^N=0$) gives $\mathcal F(\alpha_\infty)=0$.  The first
correction, at $N=4$, is $s^3v-s^4w$; therefore no coefficient of $s$ or
$s^2$ is ever changed and \eqref{eq:allorders} holds.
\end{proof}

\begin{remark}[the initial case, explicitly]\label{rem:initial}
With the lift $\gamma^{(4)}=sa_1+s^2a_2+s^3a_3$ (the coefficients of
$\bar\alpha$, $a_1=a$), the curvature is exactly
\[
 \mathcal F(\gamma^{(4)})
 =s^4\Bigl([a_1,a_3]+\tfrac12[a_2,a_2]\Bigr)
 +s^5[a_2,a_3]+s^6\tfrac12[a_3,a_3],
\]
all lower coefficients vanishing because $\bar\alpha$ is Maurer--Cartan
over $(s)/(s^4)$.  So $r_4=[a_1,a_3]+\tfrac12[a_2,a_2]$, and Step~1 for
$N=4$ asserts $\partial r_4=0$, which also follows by direct expansion.
\end{remark}

\begin{remark}[what is and is not preserved]\label{rem:preserved}
The conclusion fixes $\bar\alpha$ modulo $s^3$, \emph{not} modulo $s^4$:
the first correction $s^3v$ may alter $a_3$.  Retaining the whole
$R_4$-element would need the strictly stronger hypothesis $[r_4]=0$ in
$H^2(\h)$, which exactness does not provide and which we do not claim.
Everything downstream (Section~\ref{sec:smooth}) uses only the coefficients
of $s^0,s^1,s^2$, i.e.\ only the printed two-jet, so the weaker
preservation suffices.
\end{remark}

Combining Proposition~\ref{prop:start}, Proposition~\ref{prop:exact}, and
Lemma~\ref{lem:induction}:

\begin{corollary}\label{cor:mcinfty}
There is $\alpha_\infty\in\MC(s\,\h[[s]])$ with
$\alpha_\infty\equiv\bar\alpha\pmod{s^3}$.
\end{corollary}

\section{The algebraic family over $\QQ[[s]]$}\label{sec:family}

Put $D_\infty=d+\alpha_\infty$ on $P\otimes R$ and define
\[
 F_i:=D_\infty(e_i)\in R[x_1,\ldots,x_8]
 \qquad(1\leq i\leq16).
\tag{7.1}\label{eq:cubics}
\]

\begin{lemma}\label{lem:cubicform}
Each $F_i$ is an honest polynomial: homogeneous of degree three in
$x_1,\ldots,x_8$ with coefficients in $R=\QQ[[s]]$, and
\[
 F_i=F^{(2)}_i+s^3H_i,\qquad H_i\in R[x]_3,
\tag{7.2}\label{eq:twojetform}
\]
where $F^{(2)}_1,\ldots,F^{(2)}_{16}$ are exactly the printed frozen
two-jet cubics.
\end{lemma}

\begin{proof}
$\alpha_\infty$ preserves weight, so $\alpha_\infty(e_i)$ has weight three:
each $s$-coefficient lies in the finite-dimensional space $S_3$, making
$F_i$ an element of $R\otimes S_3=R[x]_3$ (no completion in $x$ occurs).
By Corollary~\ref{cor:mcinfty} and Computation~\ref{comp:start}(5),
$F_i\equiv F^{[3]}_i\equiv F^{(2)}_i\pmod{s^3}$.
\end{proof}

Define
\[
 B=R[x_1,\ldots,x_8]/(F_1,\ldots,F_{16}),
 \qquad
 X=\Proj B\subset\PP^7_R .
\tag{7.3}\label{eq:family}
\]

\begin{proposition}[flat projective family]\label{prop:flat}
$X\to\Spec R$ is a projective, flat morphism with special fibre
$X\times_Rk=X_0$.  Moreover every weight component $B_j$ is a free
$R$-module of rank $\dim_kA_j$; in particular the Hilbert functions of all
fibres coincide with that of $A$.
\end{proposition}

\begin{proof}
$X$ is a closed subscheme of $\PP^7_R$ cut by homogeneous polynomials, so
it is projective over $R$; no algebraization theorem is required, because
\eqref{eq:cubics} are already polynomials
(Lemma~\ref{lem:cubicform}).\footnote{If one prefers to phrase the family
formally first, the compatible flat closed subschemes
$X_n=\Proj(B/s^nB)\subset\PP^7_{R_n}$ satisfy every hypothesis of
\cite[Tag 0899]{Stacks0899}: $R$ is Noetherian and $(s)$-adically complete,
$\PP^7_R\to\Spec R$ is separated and of finite type, the squares are
cartesian closed immersions, and $X_1=X_0$ is proper over $\QQ$.  The
resulting algebraization is the closed subscheme \eqref{eq:family}, by the
uniqueness part applied to the explicit cubics.  We do not need this
route.}  The special fibre is
$\Proj\bigl(S/(F_i\bmod s)\bigr)=\Proj(S/I)=X_0$.

For each $n$, the reduction $\alpha_\infty\bmod s^n$ is a Maurer--Cartan
element of $\h_n$ whose realization is $D_\infty\bmod s^n$; by
Proposition~\ref{prop:realization},
\[
 B/s^nB=R_n[x]/(F_i\bmod s^n)=H_0(P\otimes R_n,D_\infty\bmod s^n)
\]
has all weight components free over $R_n$ of rank $\dim_kA_j$.  Fix $j$;
$B_j$ is a finitely generated $R$-module (a quotient of the finite free
module $R[x]_j$).  If $sb=0$ for some $b\in B_j$, then for every $n$ the
image of $b$ in the free $R_n$-module $B_j/s^nB_j$ is killed by $s$, hence
lies in $s^{n-1}(B_j/s^nB_j)$, so $b\in s^{n-1}B_j$; the Krull intersection
theorem for the finite module $B_j$ over the Noetherian local ring $R$
forces $b=0$.  A finitely generated torsion-free module over the discrete
valuation ring $R$ is free; its rank is
$\dim_k B_j/sB_j=\dim_kA_j$.  Hence $B=\bigoplus_jB_j$ is $R$-flat.  On the
standard chart $\{x_i\neq0\}$ the coordinate ring $(B_{x_i})_0$ is a direct
summand of the localization $B_{x_i}$ of the flat module $B$, hence flat;
the charts cover $X$, so $X\to\Spec R$ is flat.
\end{proof}

\section{The generic fibre: Cohen--Macaulay, equidimensional, smooth}
\label{sec:smooth}

\subsection{Cohen--Macaulayness and equidimensionality}

Let $B_K=B\otimes_RK=K[x]/(F_1,\ldots,F_{16})K[x]$ be the homogeneous
coordinate ring of the generic fibre and, for a field extension
$L/K$, put $B_L=B_K\otimes_KL$.

\begin{lemma}[Betti semicontinuity across the family]\label{lem:betti}
For all $i$,
$\dim_K\Tor_i^{K[x]}(B_K,K)\leq\dim_k\Tor_i^{S}(A,k)$.
\end{lemma}

\begin{proof}
Choose a graded free resolution $L_\bullet\to B$ by finitely generated free
graded $R[x]$-modules ($R[x]$ is Noetherian and $B$ is a finitely generated
graded module).  Every syzygy module
$Z_i=\ker(L_i\to L_{i-1})$ is $R$-flat: inductively, in
$0\to Z_i\to L_i\to Z_{i-1}\to0$ both right terms are $R$-flat ($Z_{-1}=B$
by Proposition~\ref{prop:flat}), so $\Tor^R_1(Z_{i-1},-)=0$ keeps the
sequence exact after any base change and $Z_i$ is flat as a kernel of a
surjection of flats onto a flat.  Consequently
$L_\bullet\otimes_Rk$ is exact except at the end, i.e.\ it is a graded free
resolution of $A$ over $S=k[x]$, and likewise $L_\bullet\otimes_RK$
resolves $B_K$ over $K[x]$.

Let $N_\bullet=L_\bullet\otimes_{R[x]}R[x]/(x_1,\ldots,x_8)$, a complex of
finitely generated free $R$-modules.  Then
$H_i(N_\bullet\otimes_Rk)=\Tor_i^S(A,k)$ and
$H_i(N_\bullet\otimes_RK)=\Tor_i^{K[x]}(B_K,K)$ by the previous paragraph.
Over the discrete valuation ring $R$, for a complex of finite free modules,
$\dim_KH_i(N\otimes K)=\operatorname{rank}_RH_i(N)$ while universal
coefficients give an injection $H_i(N)\otimes k\hookrightarrow
H_i(N\otimes k)$; hence
$\dim_KH_i(N\otimes K)\leq\dim_kH_i(N)\otimes k\leq
\dim_kH_i(N\otimes k)$.
\end{proof}

\begin{proposition}\label{prop:CM}
For every field extension $L$ of $K$, the ring $B_L$ is Cohen--Macaulay of
Krull dimension $4$, and all its minimal primes $p$ satisfy
$\dim B_L/p=4$.  Consequently the geometric generic fibre
$X_{\bar K}=\Proj B_{\bar K}$ is Cohen--Macaulay with every irreducible
component of dimension exactly $3$, and every local ring of $X_{\bar K}$ at
a closed point has Krull dimension $3$.
\end{proposition}

\begin{proof}
The Hilbert function: $\dim_L(B_L)_j=\dim_K(B_K)_j
=\operatorname{rank}_RB_j=\dim_kA_j$ (Proposition~\ref{prop:flat}).  Since
$A$ has Krull dimension $4$, its Hilbert polynomial has degree $3$, and by
Hilbert--Serre dimension theory~\cite[Ch.~4]{BH} the same holds for $B_L$:
$\dim B_L=4$.

Projective dimension: $\depth_{S_+}A=4$ (Computation~\ref{comp:special}),
so the graded Auslander--Buchsbaum formula~\cite[\S1.3]{BH} gives
$\pd_SA=8-4=4$, i.e.\ $\Tor^S_i(A,k)=0$ for $i\geq5$.  By
Lemma~\ref{lem:betti}, $\Tor^{K[x]}_i(B_K,K)=0$ for $i\geq5$, so
$\pd_{K[x]}B_K\leq4$; since $\Tor$ commutes with the flat base change
$K\to L$, also $\pd_{L[x]}B_L\leq4$.  Auslander--Buchsbaum again gives
$\depth_{\mathfrak m}B_L\geq8-4=4$ where $\mathfrak m$ is the irrelevant
maximal ideal; combined with $\dim B_L=4$ and
$\depth\leq\dim$, the local ring $(B_L)_{\mathfrak m}$ is Cohen--Macaulay
of dimension~$4$, and hence so is the graded ring $B_L$ (Cohen--Macaulayness
of a finitely generated graded algebra over a field is detected at the
irrelevant maximal ideal, and localization preserves it)~\cite[\S2.1]{BH}.

Unmixedness: a Cohen--Macaulay local ring is equidimensional---every
minimal prime $q$ of $(B_L)_{\mathfrak m}$ has
$\dim(B_L)_{\mathfrak m}/q=\dim(B_L)_{\mathfrak m}$~\cite[\S2.1]{BH}.  The
minimal primes of $B_L$ are homogeneous, hence contained in
$\mathfrak m$ and in bijection with those of $(B_L)_{\mathfrak m}$, and
graded dimensions are computed at $\mathfrak m$; so $\dim B_L/p=4$ for
every minimal prime.  No minimal prime contains $(B_L)_+$ (its quotient
would have dimension $\leq1<4$), so the irreducible components of
$\Proj B_L$ are the $\Proj(B_L/p)$, of dimension $4-1=3$.  At a closed
point of $X_{\bar K}$, the local ring has dimension equal to the maximum
dimension of components through the point, namely $3$; since all components
have dimension $3$, this holds at every closed point.
\end{proof}

\subsection{Singularity as a linear-algebra incidence}

Work on the affine chart $x_v=1$ of $\PP^7$, $1\leq v\leq8$, with
coordinates $z_1,\ldots,z_7$, and let $J$ denote the $16\times7$ Jacobian
matrix $\bigl(\partial\widetilde F_i/\partial z_u\bigr)$ of the
dehomogenized cubics.  Let $\bar K$ be an algebraic closure of $K$.

\begin{lemma}[Jacobian dichotomy]\label{lem:jacobian}
A closed point $p$ of $X_{\bar K}$ lying in the chart $x_v=1$ is a singular
point of $X_{\bar K}$ if and only if
$\rank J(p)\leq3$, if and only if $\ker J(p)\subseteq\bar K^7$ contains a
four-dimensional subspace.
\end{lemma}

\begin{proof}
The chart $X_{\bar K}\cap\{x_v\neq0\}$ is the affine scheme cut out by
exactly the sixteen dehomogenized cubics: for a graded quotient
$B_{\bar K}=\bar K[x]/(F_i)$, the degree-zero part of the localized ideal
$\bigl((F_i)_{x_v}\bigr)_0$ is generated by the $F_i/x_v^3$, since any
$g/x_v^m$ with $g=\sum h_iF_i$ homogeneous of degree $m$ equals
$\sum(h_i/x_v^{m-3})(F_i/x_v^3)$.  Hence the Zariski tangent space of
$X_{\bar K}$ at $p$ is $\ker J(p)$, of dimension $7-\rank J(p)\geq3$.  By
Proposition~\ref{prop:CM}, $\dim\mathcal O_{X_{\bar K},p}=3$.  The point is
smooth iff $\mathcal O_{X_{\bar K},p}$ is regular iff
$\dim_{\bar K}\ker J(p)=3$ iff $\rank J(p)=4$; it is singular iff
$\rank J(p)\leq3$ iff $\dim\ker J(p)\geq4$.
\end{proof}

For a four-subset $\Pi\subset\{1,\ldots,7\}$ let $K_\Pi$ be the
$7\times4$ matrix whose rows indexed by $\Pi$ form the identity and whose
remaining twelve entries are indeterminates $a_1,\ldots,a_{12}$; these are
the $35$ standard charts of $\Gr(4,7)$.  The \emph{truncated incidence
generators} on the pair $(v,\Pi)$ are the $80$ polynomials
\[
 \widetilde F^{(2)}_1,\ldots,\widetilde F^{(2)}_{16},\quad
 \bigl(J^{(2)}\bigr)\,K_\Pi
 \ \in\ \QQ[z_1,\ldots,z_7,a_1,\ldots,a_{12}][s],
\tag{8.1}\label{eq:incgens}
\]
where the superscript $(2)$ denotes the printed two-jet cubics and their
Jacobian ($J^{(2)}K_\Pi$ has $16\cdot4$ entries; the frozen scripts label
it via the transpose convention, which changes nothing).

\begin{computation}[the 280 incidence memberships]\label{comp:incidence}
For every one of the $8\cdot35=280$ pairs $(v,\Pi)$, the frozen scripts
\path{../referee-packet/code/verify_incidence_chart.m2} and
\path{../referee-packet/code/verify_all_incidence.py} prove by exact
rational module Gr\"obner reduction the membership
\[
 s^2\in\bigl(q_1,\ldots,q_{80}\bigr)+(s^3)
 \subset\QQ[z,a][s],
\tag{8.2}\label{eq:membership}
\]
where $q_1,\ldots,q_{80}$ are the generators \eqref{eq:incgens}: the $80$
generators are encoded over $\QQ[z,a][s]/(s^3)$ as a $3\times240$ module
matrix and $(0,0,1)^{\mathsf T}$ is exactly reduced to zero against it.
$277$ pairs certify at degree bound six and pairs $(2,7),(6,12),(8,16)$ at
degree bound eight; the driver checks that all $280$ pairs occur and
re-verifies one retained lift.  Replayed line: \texttt{proved 280/280 chart
pairs}.  Certificate:
\path{../referee-packet/verification/incidence_all_QQ.txt}.

\emph{What this does and does not certify}: \eqref{eq:membership} is a
family of $280$ polynomial identities about the \emph{two-jet}
\eqref{eq:incgens} only.  It makes no reference to the higher-order
coefficients $H_i$, and it is not, by itself, a smoothness statement; the
geometric content is extracted in Proposition~\ref{prop:smoothgeneric}.
\end{computation}

\subsection{Smoothness of the geometric generic fibre}

\begin{proposition}\label{prop:smoothgeneric}
The geometric generic fibre $X_{\bar K}$ is smooth.  Indeed, this holds for
every flat projective family whose equations satisfy
\eqref{eq:twojetform} with the printed two-jet.
\end{proposition}

\begin{proof}
Suppose $X_{\bar K}$ has a singular closed point.  By
Lemma~\ref{lem:jacobian}, on some chart $x_v=1$ there is a pair
$(p,\Lambda)$ of a point and a four-plane $\Lambda\subseteq\ker J(p)$.
The incidence scheme
\[
 \Sigma=\Bigl\{(p,\Lambda)\in
 (X\cap\{x_v\ne0\})\times\Gr(4,7):\ J(p)\cdot\Lambda=0\Bigr\}
\]
(over all $v$) is a scheme of finite type over $K$, and it has a
$\bar K$-point; hence it is nonempty and has a closed point whose residue
field $K'$ is a finite extension of $K$ (Hilbert Nullstellensatz for
finite-type schemes over a field).

Let $R'$ be the integral closure of $R=\QQ[[s]]$ in $K'$.  Since $R$ is a
complete discrete valuation ring, $R'$ is again a complete discrete
valuation ring, finite over $R$~\cite[Ch.~II]{SerreLF}; write $v'$ for its
valuation, so $v'(s)>0$.

The point $p$ is a $K'$-point of the projective scheme
$X$; by the valuative criterion of properness it extends to an $R'$-point
of $X\subset\PP^7_{R'}$.  Scale its homogeneous coordinates to lie in $R'$
with at least one of minimal valuation a unit; call that coordinate
$x_{v_0}$.  The chart may change ($v_0\neq v$ is possible), so re-choose
the plane there: $p$ lies in both charts, the Zariski tangent space
$T_pX_{K'}$ is intrinsic, and each chart is cut out by its sixteen
dehomogenized cubics (proof of Lemma~\ref{lem:jacobian}), so
$\dim_{K'}\ker J_{v_0}(p)=\dim_{K'}T_pX_{K'}
=\dim_{K'}\ker J_v(p)\geq4$, the last inequality because
$J_v(p)\Lambda=0$ with $\dim\Lambda=4$.  Choose a four-plane
$\Lambda'\subseteq\ker J_{v_0}(p)$ over $K'$.  As a $K'$-point of the
projective scheme $\Gr(4,7)$, $\Lambda'$ likewise extends to an
$R'$-point; scaling its Pl\"ucker coordinates to $R'$ with a
unit entry places it in a standard chart $\Pi_0$, where it is represented
by a matrix $K_{\Pi_0}$ with the twelve entries in $R'$, pivot rows the
identity, and column span reducing to $\Lambda'$ over $K'$.  The incidence
entries $J_{v_0}\,K_{\Pi_0}$, zero over $K'$ because the columns lie in
$\ker J_{v_0}(p)$, lie in $R'$ and hence vanish in $R'$, since
$R'\hookrightarrow K'$.

Fix the pair $(v_0,\Pi_0)$ and let
$(\zeta,\alpha)\in(R')^7\times(R')^{12}$ be these integral coordinates.
Membership \eqref{eq:membership}, lifted to the polynomial ring, is an
identity
\[
 s^2=\sum_{j=1}^{80}c_j\,q_j+s^3c_0
 \qquad\text{in }\QQ[z,a][s],
\tag{8.3}\label{eq:identity}
\]
with explicit polynomials $c_j,c_0$.  By \eqref{eq:twojetform}, the actual
incidence generators of the family---the dehomogenized $F_i$ and the
entries of $J_{v_0}\,K_{\Pi_0}$---differ from the truncated $q_j$ by $s^3$
times polynomials in $z,a$ with coefficients in $R$: dehomogenizing and
differentiating $s^3H_i$ with respect to a $z_u$ preserves the factor
$s^3$.  Every actual generator vanishes at $(\zeta,\alpha)$: the
dehomogenized $F_i$ vanish because the $R'$-section lies in $X$ and the
chart ideal is generated by them (proof of Lemma~\ref{lem:jacobian}), and
the entries of $J_{v_0}K_{\Pi_0}$ vanish by the previous paragraph.
Substituting $q_j=q_j'-s^3\rho_j$ into \eqref{eq:identity} and evaluating
at $(\zeta,\alpha)$ therefore yields
\[
 s^2=s^3c,\qquad c\in R'.
\]
Taking valuations: $2v'(s)=3v'(s)+v'(c)\geq3v'(s)$, contradicting
$v'(s)>0$.  Hence no singular closed point exists, and $X_{\bar K}$, being
of finite type over the algebraically closed field $\bar K$ with all closed
points smooth, is smooth.
\end{proof}

\begin{proof}[Proof of Theorem~\textup{\ref{thm:main}}]
Proposition~\ref{prop:start} converts the certified $R_4$ family into
$\bar\alpha\in\MC(\h\otimes(s)/(s^4))$ with the printed two-jet;
Proposition~\ref{prop:exact} establishes exactness of \eqref{eq:sequence}
at $H^2(\h)$ for its linear coefficient; Lemma~\ref{lem:induction} extends
$\bar\alpha$ to $\alpha_\infty\in\MC(s\,\h[[s]])$ without changing the
coefficients of $s$ and $s^2$; Lemma~\ref{lem:cubicform} and
Proposition~\ref{prop:flat} produce the flat projective family
\eqref{eq:family} with special fibre $X_0$ and equations
\eqref{eq:twojetform}; and Proposition~\ref{prop:smoothgeneric} proves its
geometric generic fibre smooth.
\end{proof}

\section{Dependency ledger}\label{sec:ledger}

\begin{center}
\begin{tabular}{p{0.22\textwidth}p{0.36\textwidth}p{0.32\textwidth}}
\toprule
Step & Status & Where\\
\midrule
Special fibre CM, $\depth=4$ & Verified computation + Reisner &
Comp.~\ref{comp:special}\\
Tate stages complete & Verified computation (M2 + Singular) &
Comp.~\ref{comp:tate}\\
Controller dimensions $109/56/53$, $27$ & Verified computation &
Comp.~\ref{comp:controller}\\
MC $\Rightarrow$ flat embedded family & Directly proved &
Prop.~\ref{prop:realization}\\
$R_4$ family $\Rightarrow$ MC element & Verified computation + directly
proved & Comp.~\ref{comp:start}, Prop.~\ref{prop:start}\\
Exactness at $H^2$ & Verified computation (ranks $12$, $15$, zero
composite, full $H^2$ basis) + directly proved transfer &
Comp.~\ref{comp:bracket}, Prop.~\ref{prop:exact}\\
All-orders MC element & Directly proved (strict Bianchi induction) &
Lem.~\ref{lem:induction}\\
Flat projective family over $R$ & Directly proved &
Prop.~\ref{prop:flat}\\
Generic fibre CM equidimensional & Directly proved (Betti
semicontinuity + Auslander--Buchsbaum) & Prop.~\ref{prop:CM}\\
Generic smoothness & Verified computation ($280$ memberships) + directly
proved valuation argument & Comp.~\ref{comp:incidence},
Prop.~\ref{prop:smoothgeneric}\\
\bottomrule
\end{tabular}
\end{center}

External classical inputs: Reisner's criterion~\cite{Reisner}; Tate's
construction (re-proved as Lemma~\ref{lem:tateexists}); local flatness
criterion \cite[Tag 00MK]{Stacks00MK} (re-proved in the case used);
Auslander--Buchsbaum, Hilbert--Serre, and unmixedness of Cohen--Macaulay
local rings~\cite{BH}; valuative criterion of properness; integral closure
of a complete DVR in a finite extension~\cite{SerreLF}.  The optional
footnote route uses \cite[Tag 0899]{Stacks0899}.  No statement about
structures transferred to cohomology, or about operations of arity
$\geq3$, is used, cited, or silently reproduced.

\begin{thebibliography}{9}

\bibitem{Avramov}
L.~L.~Avramov,
\emph{Infinite free resolutions}, in: Six Lectures on Commutative Algebra
(Bellaterra, 1996), Progress in Mathematics 166, Birkh\"auser, Basel,
1998, pp.~1--118.

\bibitem{BH}
W.~Bruns and J.~Herzog,
\emph{Cohen--Macaulay Rings}, Cambridge Studies in Advanced Mathematics 39,
Cambridge University Press, 1993.

\bibitem{GS}
B.~Gr\"unbaum and V.~P.~Sreedharan,
\emph{An enumeration of simplicial $4$-polytopes with $8$ vertices},
J.~Combinatorial Theory \textbf{2} (1967), 437--465,
\href{https://doi.org/10.1016/S0021-9800(67)80055-3}
{doi:10.1016/S0021-9800(67)80055-3}.

\bibitem{Reisner}
G.~A.~Reisner,
\emph{Cohen--Macaulay quotients of polynomial rings},
Adv.~Math. \textbf{21} (1976), 30--49,
\href{https://doi.org/10.1016/0001-8708(76)90114-6}
{doi:10.1016/0001-8708(76)90114-6}.

\bibitem{SerreLF}
J.-P.~Serre,
\emph{Local Fields}, Graduate Texts in Mathematics 67, Springer, 1979.

\bibitem{Stacks00MK}
The Stacks Project Authors, \emph{The Stacks Project},
\href{https://stacks.math.columbia.edu/tag/00MK}{Tag 00MK}
(local criterion of flatness: for a local homomorphism of Noetherian local
rings and a finite module $M$ over the target,
$\Tor_1(\kappa,M)=0$ implies $M$ flat).

\bibitem{Stacks0899}
The Stacks Project Authors, \emph{The Stacks Project},
\href{https://stacks.math.columbia.edu/tag/0899}{Tag 0899}
(algebraization of compatible systems of proper closed subschemes over a
complete Noetherian base).

\bibitem{Tate}
J.~Tate,
\emph{Homology of Noetherian rings and local rings},
Illinois J.~Math. \textbf{1} (1957), 14--27.

\end{thebibliography}

\end{document}
